\section{Safety Obligation Construction}
\label{sec:construction}

\subsection{Overview}

The purpose of obligation construction is to turn dispersed API prose and bug
evidence into a reusable checking specification. Construction occurs offline
and is independent of a target project. Each specification entry records a
trigger, the safety property induced by the trigger, the program subject, and
the conditions that discharge the property. Detection then instantiates these
entries at concrete API calls, as described in
Section~\ref{sec:pipeline}.

The specification contains only semantic information: a safety obligation and
the conditions that discharge it. Operational optimizations are deliberately
kept outside this representation. For example, the absence of a nearby
\code{Py\_DECREF} cannot define an ownership violation because the reference
may instead be transferred or returned. Section~\ref{sec:pipeline} introduces
such syntactic patterns as part of static pruning.

\subsection{Evidence Sources}

The primary source is the CPython 3.14 Python/C API reference manual
\cite{pythoncapi}. We inspect chapters that specify object ownership,
reference counting, exception handling, memory allocation, buffer access,
thread state, garbage-collection support, container iteration, argument
conversion, and capsules. Normative statements identify the trigger and the
required property. Notes and examples provide admissible safe outcomes. We
also inspect real Python/C API bug-fix commits. We mined 442 fix commits
from 23 open-source Python extension repositories via the GitHub API,
filtering to commits that modify C or C++ source files and whose messages or
diffs reference Python/C API constructs; 395 of the 442 commits involve
CPython FFI code. These cases expose requirements that are easy to overlook in
documentation, interactions between obligations, and recurring code forms
suitable for pruning patterns. We stopped adding obligation categories after
three consecutive documentation sections and ten consecutive bug cases
introduced no new obligation type. The empirical studies by Yang and Cai and
by Hu and Zhang provide independent corroboration that ownership, exception,
and compatibility mistakes are the dominant failure classes in real-world
cross-language code \cite{yang2025dissecting,hu2020pythoncapi}.

We use bug cases to refine and validate documented requirements, not to infer
correctness from frequency. A recurring idiom is not considered safe when it
contradicts the API documentation. Conversely, an unusual idiom can discharge
an obligation when its path and ownership semantics satisfy the documented
property.

\subsection{Construction Procedure}

The construction procedure has three steps. First, we extract a candidate
requirement and record its triggering APIs, subject, failure mode, and cited
documentation. For the reference-counting rule that a new reference belongs
to its caller, for example, the trigger is a new-reference-returning API, the
subject is its result, and the failure modes include leaks and premature
release.

Second, we normalize requirements by the safety property rather than by API
name. APIs that create the same duty share one obligation even when their
surface signatures differ. Requirements remain separate when they require
different evidence. Reference ownership and null validation can therefore be
bound to the same call without collapsing their meanings.

Third, we derive discharge conditions from documented safe alternatives and
validated fixes. Conditions are intentionally sufficient rather than
exhaustive. An owned reference, for example, can be released on every exit,
transferred to an ownership-consuming API, or returned to the caller. A
condition records the semantic outcome, not one exact syntax. We merge
conditions that establish the same outcome and retain separate conditions
when they require different reasoning.

\subsection{Safety Obligations}

Table~\ref{tab:obligations} summarizes the resulting specification. The 13
categories cover reference lifecycle, failure handling, native resources,
type and value agreement, and runtime protocols. A single API can bind several
categories. For example, an object-producing call can create ownership,
null-validation, cleanup, and exception-state obligations. Conversely, one
category can cover many APIs that share a contract.

\begin{table*}[t]
\caption{Python/C API safety obligation types and their discharge conditions.}
\label{tab:obligations}
\centering
\fontsize{6.2}{6.7}\selectfont
\setlength{\tabcolsep}{2.1pt}
\renewcommand{\arraystretch}{1.08}
\begin{tabularx}{\textwidth}{|>{\centering\bfseries\arraybackslash}p{2.2cm}|>{\centering\arraybackslash}p{6.5cm}|>{\centering\bfseries\arraybackslash}p{3.1cm}|X|}
\hline
\rowcolor{black!10}
\rule{0pt}{2.4ex}\textbf{Obligation type} & \textbf{Obligation description} & \textbf{Discharge condition} & \multicolumn{1}{c|}{\textbf{Discharge-condition description}} \\
\hline
\multirow[c]{4}{=}{\centering Reference ownership} &
\multirow[c]{4}{=}{\centering Resolve every owned reference and prevent reuse after ownership transfer.} &
Explicit release & Release the reference on every exit path. \\
\cline{3-4}
& & Ownership transfer & Pass it to an API that assumes release responsibility. \\
\cline{3-4}
& & Return to caller & Return it as the function's new-reference result. \\
\cline{3-4}
& & Post-steal non-use & Do not use or release it after a steal unless first retained. \\
\hline
\multirow[c]{2}{=}{\centering Reference validity} &
\multirow[c]{2}{=}{\centering Keep a borrowed reference valid at every use.} &
Scope-safe use & Finish all uses before any operation that can invalidate it. \\
\cline{3-4}
& & Strong promotion & Retain it before invalidation or persistent storage. \\
\hline
\multirow[c]{3}{=}{\centering Null/error validation} &
\multirow[c]{3}{=}{\centering Prove a fallible result valid before unsafe use.} &
Dominating check & Check NULL or the error sentinel before every unsafe use. \\
\cline{3-4}
& & Infallible source & Obtain the value from a source infallible for this input. \\
\cline{3-4}
& & NULL-safe use & Use the value only in NULL-tolerant operations. \\
\hline
\multirow[c]{4}{=}{\centering Error propagation} &
\multirow[c]{4}{=}{\centering Handle or propagate failures and clean resources on failing paths.} &
Cleanup block & Release acquired resources and return an error indicator. \\
\cline{3-4}
& & Early return & Release resources inline before returning the error. \\
\cline{3-4}
& & Recovery & Clear the exception and execute an intentional fallback. \\
\cline{3-4}
& & NULL-safe cleanup & Initialize cleanup variables and use NULL-safe releases. \\
\hline
\multirow[c]{4}{=}{\centering Exception state} &
\multirow[c]{4}{=}{\centering Resolve pending exceptions before normal API work and preserve prior state.} &
Immediate check & Inspect exception state before a non-safe API call. \\
\cline{3-4}
& & Safe-only work & Perform only exception-safe operations while one is pending. \\
\cline{3-4}
& & Prompt resolution & Clear and recover, or return an error, before normal work. \\
\cline{3-4}
& & Dealloc preservation & Save and restore pre-existing state around finalization. \\
\hline
\multirow[c]{4}{=}{\centering Resource protocol} &
\multirow[c]{4}{=}{\centering Pair each acquisition with its correct release and allocator domain.} &
Matched release & Apply the matching release on every successful-acquire path. \\
\cline{3-4}
& & Domain consistency & Free memory through the allocator domain that created it. \\
\cline{3-4}
& & Failure-aware skip & Skip release when acquisition failed and returned NULL. \\
\cline{3-4}
& & Safe reallocation & Preserve the original pointer until reallocation succeeds. \\
\hline
\multirow[c]{5}{=}{\centering Post-release safety} &
\multirow[c]{5}{=}{\centering Prevent repeated release, later access, and reentrant inconsistent state.} &
NULL assignment & Clear the pointer immediately after releasing it. \\
\cline{3-4}
& & No later access & Do not read, dereference, or pass the released resource. \\
\cline{3-4}
& & Container-safe update & Use Py\_CLEAR or Py\_SETREF for observable members. \\
\cline{3-4}
& & Local-only scope & Ensure reentrant code cannot observe the released object. \\
\cline{3-4}
& & Buffer-alias lifetime & Finish alias reads before releasing the underlying buffer. \\
\hline
\multirow[c]{2}{=}{\centering Type agreement} &
\multirow[c]{2}{=}{\centering Match API formats with C types, widths, encodings, and protocols.} &
Format-type match & Match each format code to the corresponding C argument type. \\
\cline{3-4}
& & Platform width & Use format codes consistent with platform-dependent widths. \\
\hline
\multirow[c]{3}{=}{\centering Numeric and bounds safety} &
\multirow[c]{3}{=}{\centering Prevent narrowing, allocation overflow, and out-of-bounds access.} &
Range check & Check the range before narrowing or indexed access. \\
\cline{3-4}
& & Allocation check & Check size arithmetic for overflow before allocation. \\
\cline{3-4}
& & Sufficient type & Store the value in a type wide enough for the platform. \\
\hline
\multirow[c]{2}{=}{\centering Thread safety} &
\multirow[c]{2}{=}{\centering Use the Python/C API under a valid interpreter-lock context.} &
Valid context & Hold the GIL whenever invoking the Python/C API. \\
\cline{3-4}
& & Blocking-I/O release & Release the GIL for blocking work and reacquire it first. \\
\hline
\multirow[c]{3}{=}{\centering GC lifecycle} &
\multirow[c]{3}{=}{\centering Complete allocation, tracking, traversal, clearing, and deallocation.} &
GC allocation domain & Use matching GC-aware allocation and deallocation. \\
\cline{3-4}
& & Tracking protocol & Track after initialization and untrack before field clearing. \\
\cline{3-4}
& & Traverse and clear & Visit every owned field and clear it with Py\_CLEAR. \\
\hline
\multirow[c]{2}{=}{\centering Iteration integrity} &
\multirow[c]{2}{=}{\centering Do not structurally mutate a container during iteration.} &
Read-only body & Perform no direct or indirect mutation while iterating. \\
\cline{3-4}
& & Iteration copy & Iterate over a copy when the original may be mutated. \\
\hline
\multirow[c]{3}{=}{\centering Capsule name stability} &
\multirow[c]{3}{=}{\centering Keep a capsule name valid and consistent throughout its lifetime.} &
String literal & Use an immutable string literal with static lifetime. \\
\cline{3-4}
& & Static storage & Store the name in a static or module-lifetime constant. \\
\cline{3-4}
& & Destructor-managed & Free a dynamic name only in the capsule destructor. \\
\hline
\end{tabularx}
\end{table*}

\subsection{Discharge Conditions}

Let $o$ be an obligation instance bound in function $f$, let $D(o)$ be its
discharge conditions, and let $\Pi(o,f)$ be the relevant paths from its trigger
to an exit or unsafe use. For the default disjunctive mode, the intended
judgment is
\begin{equation}
  \mathit{Safe}(o,f) \equiv
  \forall \pi \in \Pi(o,f),\; \exists d \in D(o): d(\pi).
  \label{eq:discharge}
\end{equation}
Thus, different paths can satisfy different conditions. SO-11 uses a
conjunctive mode because allocation, tracking, traversal, clearing, and
deallocation are complementary aspects of one garbage-collection protocol.
For this category, all applicable conditions must hold.

Equation~\ref{eq:discharge} defines the question asked by the checker. It does
not claim that the LLM enumerates paths soundly or proves the equation. The
structured output records the paths and condition judgments used to reach a
verdict so that a reviewer can inspect the reasoning.

The ownership obligation in Figure~\ref{fig:motivation} has four conditions:
explicit release, ownership transfer, return to the caller, and no use after a
steal. On each failure path of the new-reference branch, explicit release
holds. On the successful branch, returned arrays transfer ownership. The
borrowed-reference branch does not create an owned reference, so its cleanup
must not apply the release action. Binding ownership semantics to the concrete
assignment exposes this mismatch.

The same representation applies beyond ownership. Null validation can be
discharged by a dominating check, an infallible source for the given input, or
exclusively null-tolerant uses. A resource protocol requires matched release,
allocator-domain agreement, and a failure-safe reallocation pattern. Capsule
names can use literal, static, or destructor-managed storage. These conditions
specify semantic outcomes rather than one cleanup-label name, macro, or branch
shape, allowing equivalent C, C++, and generated-wrapper idioms.

\subsection{Traceability and Evolution}

Each obligation entry retains its trigger API categories, discharge
conditions, and documentation locations in one versioned JSON record. This
structure supports two audits. A specification audit can trace a
property or condition back to the CPython reference chapter that motivates it.
A detection audit can trace a report from a concrete call site to the bound SO,
the evaluated DCs, and the final path evidence. Keeping both links is necessary
because a correct API rule can still be bound to the wrong subject, and a
correct binding can still receive an incorrect path judgment.

The specification is designed to evolve with CPython. A new API that follows
an existing ownership or error convention usually adds a trigger without
creating a new SO. A new semantic requirement adds a condition or category
only when the existing vocabulary cannot express its safe outcomes. Static
pruning patterns are versioned separately in the analyzer, so refining a
heuristic does not change the meaning of past obligations.
